\chapter{Analyse et Spécification des besoins}

\section*{Introduction}

Dans ce chapitre, nous identifions les caractéristiques du projet en analysant les besoins
fonctionnels et non fonctionnels. Nous définissons également les acteurs et clarifions les
exigences, qui sont ensuite représentées à travers un diagramme global de cas d'utilisation.
Enfin, nous présentons l'organisation du développement en sprints.

\section{Spécification des acteurs}

\subsection{Identification des acteurs}

L'identification des acteurs permet de définir les périmètres d'accès à l'application.
Dans CrediWise, les acteurs correspondent à des profils opérationnels définis dans Keycloak
et propagés au frontend et au backend via un mécanisme d'autorisation basé sur les rôles.

Les acteurs suivants ont été identifiés :

\begin{itemize}
    \item[\ensuremath{\triangleright}] \textbf{Super Administrateur (SUPER\_ADMIN) :} dispose du niveau d'accès le plus
    élevé. Il gère les agences, les comptes gestionnaires, les rôles, les règles d'affichage
    ainsi que les données administratives du système.

    \item[\ensuremath{\triangleright}] \textbf{Utilisateur technique (TECH\_USER) :} assure des tâches d'administration
    technique, principalement liées à la gestion des agences et des comptes gestionnaires,
    sans intervenir dans les processus métier de décision.

    \item[\ensuremath{\triangleright}] \textbf{Chargé de relation client (CRO) :} consulte les dossiers clients, suit
    les demandes de crédit et complète les dossiers d'analyse lorsqu'ils lui sont affectés.

    \item[\ensuremath{\triangleright}] \textbf{Agent front office (FRONT\_OFFICE) :} assure l'accueil des clients, crée
    les fiches clients, prépare les demandes de crédit et soumet les brouillons pour
    traitement.

    \item[\ensuremath{\triangleright}] \textbf{Décideur agence (BRANCH\_DM) :} examine les demandes de crédit et les
    dossiers d'analyse au niveau de l'agence. Il participe aux décisions et aux transitions
    de statut selon les règles d'autorisation.

    \item[\ensuremath{\triangleright}] \textbf{Décideur siège (HEAD\_OFFICE\_DM) :} consulte les dossiers nécessitant
    une validation centrale et participe au processus de décision pour les demandes hors
    périmètre agence.

    \item[\ensuremath{\triangleright}] \textbf{Analyste risque (RISK\_ANALYST) :} analyse les demandes de crédit, les
    données d'évaluation et les résultats de scoring afin de contribuer à l'analyse des
    risques.

    \item[\ensuremath{\triangleright}] \textbf{Auditeur (READ\_ONLY) :} dispose d'un accès en lecture seule, utilisé
    pour les activités d'audit, de contrôle et de vérification.

    \item[\ensuremath{\triangleright}] \textbf{Fournisseur d'identité (Keycloak) :} acteur système externe chargé de
    l'authentification, de la gestion des sessions et de l'émission des jetons. Il fournit
    également les rôles consommés par le frontend et le backend.
\end{itemize}

\section{Spécification des besoins}

Dans cette section, nous analysons les besoins fonctionnels et non fonctionnels du système,
constituant une étape essentielle dans la définition de l'architecture du projet.

\subsection{Besoins fonctionnels}

\subsubsection*{\ensuremath{\triangleright}~Authentification et sécurité}

\begin{itemize}
    \item Le système authentifie les utilisateurs via Keycloak et récupère les jetons
    d'accès à l'aide de l'adaptateur keycloak-js.
    \item Le frontend protège les routes en fonction des rôles de l'utilisateur, tandis
    que le backend applique les autorisations via l'annotation \texttt{@RolesAllowed}.
    \item Les jetons d'accès sont attachés aux requêtes API internes et renouvelés
    automatiquement avant expiration afin d'assurer la continuité des sessions.
    \item La déconnexion est gérée par Keycloak, garantissant la fermeture complète de
    la session au niveau du fournisseur d'identité.
\end{itemize}

\subsubsection*{\ensuremath{\triangleright}~Gestion des utilisateurs et des agences}

\begin{itemize}
    \item Le service Manager permet de créer, modifier, désactiver et consulter les
    comptes gestionnaires.
    \item Le système prend en charge la gestion complète des agences (création, mise à
    jour, consultation et activation/désactivation).
    \item Les gestionnaires peuvent être associés à des agences, avec une cohérence
    assurée entre les rôles Keycloak et les données applicatives.
    \item Chaque utilisateur peut consulter et mettre à jour son profil selon les droits
    définis par l'API.
\end{itemize}

\subsubsection*{\ensuremath{\triangleright}~Gestion des clients}

\begin{itemize}
    \item Le service Customer gère les fiches clients (personnes physiques et morales).
    \item Les données incluent l'identité, les coordonnées, l'agence, le type de compte,
    le segment, le secteur d'activité, le scoring et le cycle client.
    \item Le système permet la recherche de clients par identifiant ou coordonnées, ainsi
    que la consultation paginée des listes.
    \item Les données de référence (segments, secteurs, activités, types de compte, etc.)
    sont exposées pour garantir une saisie structurée.
\end{itemize}

\subsubsection*{\ensuremath{\triangleright}~Gestion des demandes de crédit}

\begin{itemize}
    \item Le service Credit Request permet la création et la mise à jour des demandes de
    crédit au statut \texttt{DRAFT}.
    \item Une demande inclut les informations client, le montant, la durée, le produit,
    les garanties, les garants, la capacité de remboursement et les indicateurs de risque.
    \item Une fois soumise, la demande est routée selon le produit : les produits 101,
    102 et 103 sont dirigés vers l'analyse, les autres vers la checklist pré-comité.
    \item Les acteurs autorisés peuvent mettre à jour le statut dans le cadre du cycle
    de vie du crédit.
    \item Une version PDF de la demande peut être générée pour archivage et documentation.
\end{itemize}

\subsubsection*{\ensuremath{\triangleright}~Analyse du crédit et scoring}

\begin{itemize}
    \item Le service Analysis crée un dossier d'analyse pour chaque demande éligible.
    \item Le processus suit plusieurs étapes : informations client, objet du crédit,
    risques client, commerciaux et financiers, et données complémentaires.
    \item Les opérations de sauvegarde et de validation sont limitées aux profils autorisés.
    \item Un module de scoring calcule et enregistre le score de crédit associé à la demande.
    \item Les règles d'affichage sont versionnées afin de garantir la traçabilité des décisions.
\end{itemize}

\subsubsection*{\ensuremath{\triangleright}~Liste de contrôle pré-comité}

\begin{itemize}
    \item Le service Checklist consolide les documents et vérifications nécessaires avant
    le passage en comité.
    \item La checklist est générée ou mise à jour à partir de la demande et enrichie via
    le service Customer.
    \item Elle couvre les documents, garanties, archivage, préparation au décaissement
    et contrôles front office/back office.
\end{itemize}

\subsubsection*{\ensuremath{\triangleright}~Analyse du risque crédit}

\begin{itemize}
    \item Le service Risk Analysis gère le formulaire d'analyse du risque crédit (CRA).
    \item Le formulaire suit plusieurs statuts contrôlés : brouillon, en révision, validé,
    envoyé au comité, ou archivé.
    \item Les analystes peuvent compléter les sections d'avis et de questions-réponses.
    \item Un module d'audit enregistre l'historique des modifications.
\end{itemize}

\subsubsection*{\ensuremath{\triangleright}~Décision du comité}

\begin{itemize}
    \item Le service Credit Committee enregistre et expose les décisions du comité de crédit.
    \item Les décisions sont initialisées à partir des services amont et mises à jour après
    validation.
    \item Elles influencent directement le cycle de vie de la demande (approbation, rejet
    ou décaissement).
\end{itemize}

\subsubsection*{\ensuremath{\triangleright}~Décaissement}

\begin{itemize}
    \item Le service Disbursement enregistre les informations de décaissement pour les
    demandes approuvées.
    \item L'enregistrement est basé sur l'instantané de la demande de crédit.
    \item Une fois validé, le statut passe à « décaissé », et le compteur client est mis
    à jour via gRPC.
\end{itemize}

\subsubsection*{\ensuremath{\triangleright}~Auditabilité et traçabilité}

\begin{itemize}
    \item Chaque service enregistre les métadonnées de création et de mise à jour des entités.
    \item La sécurité contrôle les accès aux routes API et aux services backend.
    \item Les utilisateurs en lecture seule peuvent consulter les données sans pouvoir les
    modifier, garantissant les besoins d'audit et de conformité.
\end{itemize}

\subsection{Besoins non fonctionnels}

\subsubsection*{\ensuremath{\triangleright}~Performance}

Le backend est implémenté avec Quarkus, offrant un temps de démarrage rapide et une faible
consommation mémoire. Les API REST intègrent des mécanismes de pagination et de filtrage
afin de limiter la taille des réponses. Le frontend, développé avec Vite et React, assure
une navigation fluide et réactive en réduisant les rechargements de page.

\subsubsection*{\ensuremath{\triangleright}~Scalabilité}

L'architecture repose sur une séparation des principaux domaines métier en services
indépendants : Manager, Customer, Credit Request et Analysis. Cette approche permet à chaque
service, ainsi qu'à sa base de données associée, d'évoluer et d'être dimensionné
indépendamment en fonction de la charge.

\subsubsection*{\ensuremath{\triangleright}~Fiabilité et disponibilité}

L'infrastructure locale est déployée via Docker Compose, incluant PostgreSQL, Keycloak et
pgAdmin. La séparation des services et des bases de données limite l'impact des défaillances
à un domaine spécifique et garantit la reproductibilité des environnements de développement.

\subsubsection*{\ensuremath{\triangleright}~Maintenabilité}

Chaque service backend est structuré comme un projet Maven indépendant avec son propre code,
sa configuration, ses tests et ses migrations Flyway. Cette organisation réduit le couplage
entre services et facilite la maintenance ainsi que l'évolution du système.

\subsubsection*{\ensuremath{\triangleright}~Conformité}

Le système traite des données personnelles et financières, ce qui impose des exigences
strictes de traçabilité et de contrôle d'accès. L'utilisation d'instantanés immuables des
demandes et de métadonnées d'audit garantit la cohérence des données et la conformité aux
exigences réglementaires.

\subsubsection*{\ensuremath{\triangleright}~Sécurité}

L'authentification est centralisée via Keycloak. Les services backend valident les jetons
JWT et extraient les rôles configurés pour appliquer les règles d'autorisation. Les contrôles
d'accès sont appliqués au niveau des API et renforcés au niveau métier lorsque nécessaire.
Les jetons ne sont pas stockés dans le stockage local du navigateur afin de réduire les
risques de sécurité.

\subsubsection*{\ensuremath{\triangleright}~Ergonomie et expérience utilisateur}

Le frontend est une application monopage (SPA) développée avec React et TypeScript. Il
propose une navigation basée sur les rôles, une gestion multilingue via i18next, une
validation des formulaires et un système cohérent de gestion des erreurs afin d'améliorer
l'expérience utilisateur.

\subsection{Backlog de Produit}

\setcounter{table}{0}
\renewcommand{\thetable}{1.\arabic{table}}

Après avoir défini les besoins fonctionnels de notre solution, le backlog de produit est présenté
dans le tableau~\ref{tab:backlog} ci-après :

\begin{longtable}{|p{3cm}|p{8.5cm}|c|c|}
\hline
\textbf{Sprint} &
\textbf{UserStory} &
\textbf{Priorité} &
\textbf{Estimation} \\
\hline
\endfirsthead
\hline
\textbf{Sprint} &
\textbf{UserStory} &
\textbf{Priorité} &
\textbf{Estimation} \\
\hline
\endhead
\endlastfoot

% -------- Sprint 1 --------
\multirow{10}{3cm}{\textbf{Sprint 1 : Authentification et gestion de la clientèle}}
& En tant que gestionnaire, je veux me connecter à la plateforme avec mes identifiants Keycloak, afin d'accéder à mon espace de travail sécurisé. & 1 & 8h \\ \cline{2-4}
& En tant que gestionnaire, je veux que mon rôle détermine automatiquement mes droits d'accès. & 1 & 10h \\ \cline{2-4}
& En tant que super administrateur, je veux créer et gérer les comptes gestionnaires .& 1 & 12h \\ \cline{2-4}
& En tant que super administrateur, je veux créer et gérer les agences. & 1 & 8h \\ \cline{2-4}
& En tant que gestionnaire, je veux consulter et modifier mon profil, afin de maintenir mes données à jour. & 2 & 5h \\ \cline{2-4}
& En tant que gestionnaire, je veux créer et gérer un client  & 1 & 12h \\ \cline{2-4}
& En tant que gestionnaire, je veux rechercher et filtrer les clients par nom, CIN ou segment, afin de retrouver rapidement un dossier existant. & 2 & 6h \\ \cline{2-4}
& En tant que super administrateur, je veux suspendre ou réactiver un compte gestionnaire, afin d'appliquer les politiques RH de l'établissement. & 3 & 4h \\ \hline

% -------- Sprint 2 --------
\multirow{10}{3cm}{\textbf{Sprint 2 : Service de Demande de Crédit}}
& En tant que gestionnaire \texttt{FRONT\_OFFICE}, je veux créer une demande de crédit en brouillon pour un client (montant, durée, produit, objet), afin d'initier le processus d'octroi. & 1 & 12h \\ \cline{2-4}
& En tant que gestionnaire, je veux prendre un instantané des données client au moment de la soumission, afin de conserver une trace immuable des informations utilisées. & 1 & 8h \\ \cline{2-4}
& En tant que gestionnaire, je veux ajouter des garants et des garanties (collatéral, hypothèque) à la demande, afin de sécuriser le dossier de crédit. & 1 & 10h \\ \cline{2-4}
& En tant que gestionnaire, je veux soumettre la demande de crédit, afin de déclencher automatiquement son acheminement vers le bon service selon le produit choisi. & 1 & 8h \\ \cline{2-4}
& En tant que gestionnaire, je veux modifier une demande en statut brouillon, afin de corriger les informations avant soumission. & 2 & 5h \\ \cline{2-4}
& En tant que gestionnaire, je veux consulter la liste de toutes les demandes avec leur statut (brouillon, soumis, en analyse, décaissé, rejeté), afin de suivre l'avancement du portefeuille. & 2 & 6h \\ \cline{2-4}
& En tant que gestionnaire, je veux supprimer une demande en statut brouillon, afin d'annuler une saisie erronée. & 3 & 3h \\ \cline{2-4}
& En tant que système, je veux router automatiquement les produits 101/102/103 vers le dossier d'analyse et les produits 104/105 vers la checklist pré-comité, afin de respecter le circuit d'approbation propre à chaque produit. & 1 & 8h \\ \cline{2-4}
& En tant que gestionnaire, je veux accéder au catalogue des produits de crédit disponibles, afin de choisir le bon produit lors de la création d'une demande. & 2 & 4h \\ \cline{2-4}
& En tant que gestionnaire, je veux voir le score de risque calculé pour la demande soumise, afin d'avoir une première évaluation du niveau de risque. & 3 & 6h \\ \hline

% -------- Sprint 3 --------
\multirow{11}{3cm}{\textbf{Sprint 3 : Service d'Analyse}}
& En tant qu'analyste de crédit, je veux accéder au dossier d'analyse en 7 étapes associé à une demande soumise, afin de conduire l'instruction du dossier de manière structurée. & 1 & 10h \\ \cline{2-4}
& En tant qu'analyste, je veux renseigner l'étape 1 (visite client) avec la date de visite, le lieu et les observations terrain, afin de documenter le contact avec le client. & 1 & 8h \\ \cline{2-4}
& En tant qu'analyste, je veux renseigner l'étape 2 (objet du crédit) avec le type d'actif et la finalité du financement, afin de qualifier l'investissement. & 1 & 6h \\ \cline{2-4}
& En tant qu'analyste, je veux évaluer le risque client à l'étape 3 (incidents bancaires, litiges, restrictions), afin d'identifier les signaux d'alerte. & 1 & 8h \\ \cline{2-4}
& En tant qu'analyste, je veux évaluer le risque commercial à l'étape 4 (marché, concurrence, environnement), afin d'apprécier la viabilité économique du projet. & 2 & 8h \\ \cline{2-4}
& En tant qu'analyste, je veux analyser les états financiers (bilan, compte de résultat) à l'étape 5 via extraction OCR automatique, afin d'accélérer la saisie des données financières. & 1 & 16h \\ \cline{2-4}
& En tant qu'analyste, je veux évaluer les garanties proposées à l'étape 6 (valeur, type, rang), afin de quantifier la couverture du risque. & 2 & 8h \\ \cline{2-4}
& En tant qu'analyste, je veux rédiger la proposition finale à l'étape 7 avec la recommandation et les conditions suggérées, afin de soumettre un avis motivé au comité. & 1 & 8h \\ \cline{2-4}
& En tant qu'analyste, je veux sauvegarder chaque étape en mode brouillon ou la confirmer définitivement, afin de pouvoir reprendre le travail en plusieurs sessions. & 2 & 6h \\ \cline{2-4}
& En tant qu'analyste, je veux voir la recommandation IA (APPROUVER/REJETER/RÉVISER) générée après confirmation de l'étape 7, afin d'obtenir un avis complémentaire non contraignant. & 2 & 10h \\ \cline{2-4}
& En tant qu'administrateur, je veux gérer les règles d'affichage des étapes d'analyse (actif/inactif, version), afin d'adapter le formulaire aux évolutions réglementaires. & 3 & 6h \\ \hline

% -------- Sprint 4 --------
\multirow{7}{3cm}{\textbf{Sprint 4 : Service Checklist}}
& En tant que gestionnaire, je veux générer automatiquement la checklist associée à une demande de crédit, afin de disposer d'une liste structurée des documents et validations requis. & 1 & 8h \\ \cline{2-4}
& En tant que gestionnaire, je veux consulter et cocher les éléments de la checklist pré-comité section par section, afin de m'assurer que le dossier est complet avant présentation. & 1 & 10h \\ \cline{2-4}
& En tant que gestionnaire, je veux valider les éléments de la checklist de décaissement (documents signés, conditions levées), afin de confirmer que toutes les conditions préalables sont remplies. & 1 & 8h \\ \cline{2-4}
& En tant que gestionnaire, je veux gérer les entrées de la checklist de garanties (type, valeur, statut de validation), afin de vérifier la conformité des sûretés. & 2 & 8h \\ \cline{2-4}
& En tant que gestionnaire, je veux valider les sections de la checklist SFO et contrôle, afin de respecter les procédures internes de vérification. & 2 & 6h \\ \cline{2-4}
& En tant que système, je veux déclencher une recommandation IA après la validation d'une section de checklist, afin d'alimenter en continu l'analyse assistée. & 3 & 6h \\ \cline{2-4}
& En tant que gestionnaire, je veux régénérer (rafraîchir) la checklist d'une demande, afin de réinitialiser les éléments en cas de modification du dossier. & 3 & 4h \\ \hline

% -------- Sprint 5 --------
\multirow{8}{3cm}{\textbf{Sprint 5 : Service d'Analyse des Risques}}
& En tant qu'analyste risque (\texttt{CRO}/\texttt{RISK\_ANALYST}), je veux créer un formulaire CRA (Analyse du Risque de Crédit) pour une demande, afin de formaliser l'évaluation du risque de manière indépendante. & 1 & 10h \\ \cline{2-4}
& En tant qu'analyste risque, je veux renseigner les éléments de checklist du formulaire CRA par section, afin de valider la conformité documentaire et procédurale. & 1 & 8h \\ \cline{2-4}
& En tant qu'analyste risque, je veux répondre aux questions d'assurance qualité (Q\&A) du formulaire CRA, afin de documenter les vérifications effectuées. & 1 & 8h \\ \cline{2-4}
& En tant qu'analyste risque, je veux saisir des opinions de risque par thème (avec couleur et note), afin d'émettre une appréciation qualitative structurée sur chaque dimension du risque. & 1 & 10h \\ \cline{2-4}
& En tant qu'analyste risque, je veux faire progresser le statut du formulaire CRA (\texttt{BROUILLON} $\rightarrow$ \texttt{EN\_RÉVISION} $\rightarrow$ \texttt{EN\_ATTENTE\_CRO} $\rightarrow$ \texttt{VALIDÉ} $\rightarrow$ \texttt{ENVOYÉ\_AU\_COMITÉ}), afin de piloter l'avancement du processus d'évaluation. & 2 & 6h \\ \cline{2-4}
& En tant qu'analyste risque, je veux consulter la piste d'audit du formulaire CRA (qui a fait quoi et quand), afin d'assurer la traçabilité des actions. & 2 & 6h \\ \cline{2-4}
& En tant que système, je veux déclencher une recommandation IA après la soumission d'une opinion de risque, afin d'enrichir l'évaluation par un avis algorithmique. & 3 & 6h \\ \cline{2-4}
& En tant qu'analyste risque, je veux accéder aux métadonnées de référence du formulaire CRA (sections, thèmes d'opinion, types de Q\&A), afin de bénéficier d'un formulaire normalisé. & 2 & 4h \\ \hline

% -------- Sprint 6 --------
\multirow{9}{3cm}{\textbf{Sprint 6 : Services du comité de crédit et de décaissement}}
& En tant que décideur (\texttt{BRANCH\_DM}/\texttt{HEAD\_OFFICE\_DM}), je veux accéder au dossier du comité de crédit pour une demande, afin de préparer la session de décision. & 1 & 8h \\ \cline{2-4}
& En tant que décideur, je veux enregistrer la décision du comité (APPROUVÉ, REJETÉ, CONDITIONNEL) avec le montant et la durée approuvés, afin de formaliser l'issue de la délibération. & 1 & 10h \\ \cline{2-4}
& En tant que décideur, je veux saisir les conditions d'approbation et les observations du comité, afin de préciser les engagements attendus du client. & 1 & 8h \\ \cline{2-4}
& En tant que décideur, je veux enregistrer une dérogation (motif, commentaire, valideur) lorsque la décision déroge aux politiques standards, afin d'assurer la traçabilité des exceptions. & 2 & 8h \\ \cline{2-4}
& En tant que gestionnaire du comité, je veux renseigner le formulaire de préparation (nom du rapporteur, nom du décideur), afin d'organiser la présentation du dossier. & 2 & 5h \\ \cline{2-4}
& En tant que système, je veux déclencher une recommandation IA à chaque étape clé du comité (décision, préparation, dérogation), afin d'assister le décideur par un avis objectif. & 2 & 8h \\ \cline{2-4}
& En tant que gestionnaire \texttt{FRONT\_OFFICE}, je veux créer le dossier de décaissement pour une demande approuvée, afin d'initier la mise en place du crédit. & 1 & 8h \\ \cline{2-4}
& En tant que gestionnaire, je veux enregistrer les informations de décaissement (montant, date, référence Amplitude), afin de tracer la mise à disposition des fonds. & 1 & 8h \\ \cline{2-4}
& En tant que gestionnaire, je veux confirmer le décaissement effectif, afin de clôturer le cycle d'octroi et d'incrémenter le cycle crédit du client. & 1 & 6h \\ \hline

\caption{Backlog de Produit}
\label{tab:backlog}
\end{longtable}

\subsection{Planification des sprints}

Le système sera développé de manière itérative selon la méthodologie Scrum. Chaque sprint
permettra de livrer progressivement de nouvelles fonctionnalités, contribuant à la construction
du produit final. Le tableau~\ref{tab:sprints} présente la planification des sprints :

\begin{longtable}{|c|p{8cm}|c|c|}
\hline
\textbf{Sprint} &
\textbf{Libellé} &
\textbf{Début prévu} &
\textbf{Fin prévu} \\
\hline
\endfirsthead
\hline
\textbf{Sprint} &
\textbf{Libellé} &
\textbf{Début prévu} &
\textbf{Fin prévu} \\
\hline
\endhead
\endlastfoot

\multirow{2}{*}{\textbf{1}}
& Conception générale et mise en place de l'infrastructure & 16-02-2025 & 21-02-2025 \\ \cline{2-4}
& Authentification Keycloak et gestion de la clientèle     & 24-02-2025 & 07-03-2025 \\ \hline

\multirow{2}{*}{\textbf{2}}
& Modélisation et création des demandes de crédit          & 10-03-2025 & 19-03-2025 \\ \cline{2-4}
& Soumission, scoring et routage des demandes              & 20-03-2025 & 28-03-2025 \\ \hline

\multirow{2}{*}{\textbf{3}}
& Dossier d'analyse — étapes 1 à 4 (visite client, objet, risques) & 31-03-2025 & 11-04-2025 \\ \cline{2-4}
& Analyse financière OCR et proposition finale — étapes 5 à 7      & 14-04-2025 & 25-04-2025 \\ \hline

\multirow{2}{*}{\textbf{4}}
& Checklist pré-comité et checklist de décaissement        & 28-04-2025 & 02-05-2025 \\ \cline{2-4}
& Checklist garanties, SFO et contrôles                   & 05-05-2025 & 09-05-2025 \\ \hline

\multirow{2}{*}{\textbf{5}}
& Formulaire CRA — checklist risque et questions Q\&A      & 12-05-2025 & 16-05-2025 \\ \cline{2-4}
& Opinions de risque, workflow CRA et piste d'audit        & 19-05-2025 & 23-05-2025 \\ \hline

\multirow{2}{*}{\textbf{6}}
& Comité de crédit — décision, dérogation et préparation   & 26-05-2025 & 06-06-2025 \\ \cline{2-4}
& Décaissement, clôture du dossier et recommandation IA    & 09-06-2025 & 16-06-2025 \\ \hline

\caption{Planification des sprints}
\label{tab:sprints}
\end{longtable}

\section*{Conclusion}
